% !TEX root = ./mid.tex

\documentclass[11pt]{article}

% Change "final" to "review" to add line numbers for reviewing.
\usepackage[final]{acl}

% Standard package includes
\usepackage{times}
\usepackage{latexsym}
\usepackage[T1]{fontenc}
\usepackage[utf8]{inputenc}
\usepackage{microtype}
\usepackage{inconsolata}
\usepackage{graphicx}
\usepackage{hyperref}
\usepackage{booktabs}
\usepackage{multirow}
\usepackage{amsmath}
\usepackage{amssymb}

\title{Prior Case Retrieval via Rhetorical Role Segmentation\\Mid report}

\author{
  Bibek Singh Dhody \\ 2023101054 \And
  Gautam Bhetanabhotla \\ 2023101032 \And
  Raghav Grover \\ 2023101102 \And
  Sanyam Agrawal \\ 2023101038
}

\begin{document}
\maketitle

\begin{abstract}
As part of the IL-TUR (Indian legal test understanding and reasoning) benchmark for legal models, one such task is PCR (Prior Case Retrieval), where the goal is to retrieve relevant past cases given a query case document.
The state-of-the-art approach for this task is U-CREAT \cite{joshi2023ucreat}, which uses event extraction to condense factual narratives and a citation graph to propagate relevance signals.
Legal judgments are long and structurally heterogeneous, mixing facts, arguments, statutes, and procedural details.
Traditional retrieval methods such as BM25 rely on lexical matching and are unable to distinguish legally salient content from noise.
We propose a multi-stage retrieval pipeline that combines classical IR techniques, semantic embedding retrieval, and transformer-based rhetorical role segmentation to filter documents down to their most relevant sections---specifically \emph{Facts} and \emph{Ratio Decidendi}---before retrieval.
In this mid-semester report, we present extensive baseline experiments covering BM25, TF-IDF, LSA, concept extraction, graph-ensemble, and graph-diffusion methods on the Indian Legal Prior Case Retrieval (IL-PCR) benchmark.
Additionally, we demonstrate a working rhetorical role prediction system and a legal statute identification model, both of which form components of our proposed pipeline.
\end{abstract}

%----------------------------------------------------------------------
\section{Introduction}
%----------------------------------------------------------------------

The Indian legal system faces an enormous backlog of over 40 million pending cases across the Supreme Court, High Courts, and subordinate courts, making the efficient retrieval of relevant precedents a practical necessity for legal professionals.
When framing a legal document, judges and lawyers rely on their expertise to cite previous cases that support their arguments and reasoning.
However, with an exponentially growing number of cases, it becomes practically impossible to recall all relevant precedents manually.
Automating this process can make the jobs of legal experts easier and expedite the justice delivery process.

Prior Case Retrieval (PCR) is the task of ranking a corpus of candidate legal documents by relevance to a given query case.
Crucially, relevance in the legal domain differs from standard information retrieval: rather than measuring semantic similarity at the surface level, legal relevance is primarily about shared factual situations and applicable legal precedents.
PCR therefore requires a model to reason about legal documents by following the chain of facts and arguments presented, much like a judge.
Unlike standard document retrieval, legal PCR is challenging because: (i)~judgments are extremely long (averaging 4,000 tokens), (ii)~they contain structurally diverse sections with varying degrees of relevance, and (iii)~precedent similarity depends on shared legal reasoning rather than shared vocabulary.

The IL-TUR (Indian Legal Text Understanding and Reasoning) benchmark \cite{iltur2024} was introduced to systematically evaluate models on eight tasks spanning the Indian legal domain, including Legal Named Entity Recognition, Rhetorical Role Prediction, Court Judgment Prediction, Bail Prediction, Legal Statute Identification, Prior Case Retrieval, Summarization, and Legal Machine Translation.
PCR is one of the most challenging tasks in this benchmark, with even the state-of-the-art achieving only a micro-F1 of 39.15---far below the ideal baseline of 64.87.

Traditional IR methods such as BM25 treat documents as flat bags of words and struggle to distinguish legally salient content from procedural boilerplate.
Transformer-based semantic models offer richer representations but are limited by fixed context windows (typically 512 tokens), which is far shorter than a typical Indian court judgment.
Surprisingly, transformer-based models have been observed to underperform classical methods on PCR, suggesting that naive application of pretrained encoders to long legal documents is insufficient.

To address these challenges, we propose a multi-stage retrieval pipeline that:
\begin{enumerate}
    \item Applies \textbf{rhetorical role segmentation} to decompose judgments into semantically coherent units (Facts, Arguments, Ratio, Statutes, Rulings).
    \item Filters documents to retain only the most relevant sections (\emph{Facts} and \emph{Ratio Decidendi}) for retrieval.
    \item Combines classical IR baselines (BM25, TF-IDF) with semantic re-ranking using sentence-transformer embeddings.
\end{enumerate}

Our dataset and benchmark tasks are derived from the IL-PCR (Indian Legal Prior Case Retrieval) corpus of 7,070 documents \cite{joshi2023ucreat} and the IL-TUR shared task \cite{iltur2024}.
In this mid-semester report, we present the results of all completed baseline experiments, describe the rhetorical role prediction and statute identification components, and outline the remaining work.

%----------------------------------------------------------------------
\section{Related Work}
%----------------------------------------------------------------------

\paragraph{Lexical and Semantic Retrieval.}
BM25 \cite{robertson2009bm25} remains a strong baseline in legal retrieval due to the distinctive vocabulary of legal texts.
The BM25 scoring function extends TF-IDF with term frequency saturation (controlled by $k_1$) and document length normalization (controlled by $b$), which are designed to prevent long documents from being unfairly boosted.
TF-IDF with cosine similarity offers competitive performance, especially when combined with bigram features that capture short legal phrases such as ``prima facie'' or ``Section 302.''
Latent Semantic Analysis (LSA) captures latent topic structure via truncated SVD on the term-document matrix, projecting documents into a lower-dimensional semantic space where synonymous terms are mapped to similar representations.
However, its gains over direct TF-IDF are modest in legal corpora, likely because legal vocabulary is already highly specific and exhibits less synonymy than general-domain text.
Transformer-based models such as LegalBERT and InLegalBERT provide domain-specific representations through continued pretraining on legal corpora, but are constrained by the 512-token input limit, requiring chunking or hierarchical strategies for long documents.

\paragraph{Rhetorical Role Labeling.}
Rhetorical role labeling (RRL) is the task of classifying each sentence in a legal judgment into a functional category such as \emph{Fact}, \emph{Argument}, \emph{Ratio}, \emph{Ruling}, or \emph{Precedent}.
The motivation for RRL stems from the observation that legal documents, while long and unstructured, contain topically coherent units whose boundaries are predictable from discourse structure.
Prior work \cite{malik2022rr} has used multi-task BiLSTM-CRF architectures with BERT embeddings to achieve coherent label sequences across documents.
The key insight is that rhetorical roles exhibit strong \emph{label inertia}---adjacent sentences tend to share the same role, and transitions between roles are smooth rather than abrupt.
This makes sequence-labeling models (which model inter-label dependencies) substantially more effective than independent sentence classifiers.
RRL has been applied in downstream tasks including judgment summarization, argument mining, and court judgment prediction, but its use for prior case retrieval remains underexplored.

\paragraph{Event-Based and Graph Methods.}
Citation graph approaches exploit the network structure of legal corpora, where cases cite prior cases and reference statutes.
U-CREAT \cite{joshi2023ucreat} defined events as tuples of a predicate (typically a verb) and its arguments (subject, object, etc.), obtained by parsing legal documents with a dependency parser.
By representing documents through their event tuples and computing relevance via event overlap, U-CREAT achieved the current state-of-the-art micro-F1 of 39.15 on IL-PCR.
Notably, BM25 with n-gram indexing over event-filtered documents was the key driver of this performance: the best variant uses quad-gram BM25, suggesting that multi-word event expressions carry strong retrieval signal.
Our work is inspired by U-CREAT but replaces event extraction with rhetorical role segmentation, which provides a more general and interpretable decomposition of legal documents and does not require dependency parsing infrastructure.

%----------------------------------------------------------------------
\section{Dataset}
%----------------------------------------------------------------------

We use the \textbf{IL-PCR} (Indian Legal Prior Case Retrieval) corpus \cite{joshi2023ucreat}, which consists of 7,070 Indian Supreme Court judgments scraped from the publicly available IndianKanoon\footnote{\url{https://indiankanoon.org/}} portal.
All individuals' and organizations' names in the documents have been anonymized to remove potential biases in model training and evaluation.
The corpus is split into a training set (\texttt{ik\_train}) and a test set (\texttt{ik\_test}), each containing two subdirectories:
\begin{itemize}
    \item \textbf{Query cases} (\texttt{query/}): The set of documents for which relevant precedents must be retrieved.
    \item \textbf{Candidate cases} (\texttt{candidate/}): The corpus of documents to be ranked by relevance to a given query.
\end{itemize}

Table~\ref{tab:dataset} summarises the dataset statistics.
Each document is a plain-text file containing the full judgment text, including structured header fields (petitioner, respondent, date of judgment, bench, citations, act, headnote) followed by the judgment body.
Ground-truth relevance labels are provided as a JSON file mapping each query document to its list of relevant candidate filenames.
Each query has a variable number of relevant candidates (ranging from 1 to over a dozen).

\begin{table}[t]
\centering
\small
\begin{tabular}{@{}lrr@{}}
\toprule
& \textbf{Train} & \textbf{Test} \\
\midrule
Query cases     & 827  & 237  \\
Candidate cases & 4320 & 1727 \\
\bottomrule
\end{tabular}
\caption{IL-PCR dataset statistics.}
\label{tab:dataset}
\end{table}

Evaluation is performed at cutoffs $K \in \{1, 5, 10, 20\}$ using standard IR metrics: Precision@$K$, Recall@$K$, Micro-F1, Macro-F1, MAP, MRR, and NDCG@$K$.
The official IL-TUR evaluation metric for PCR is \textbf{micro-averaged F1@$K$}.
Each candidate is assigned a relevance score by the model; the top-$K$ candidates are predicted as relevant, and micro-F1 is computed against the ground-truth labels.
The choice of micro-averaging reflects the fact that queries have highly variable numbers of relevant candidates (from 1 to over a dozen), and micro-averaging gives equal weight to each query--candidate pair rather than each query.

%----------------------------------------------------------------------
\section{Baseline Experiments}
%----------------------------------------------------------------------

We conducted an extensive set of baseline experiments organized into six experimental blocks, progressively increasing the sophistication of document representation from raw lexical features to latent semantic spaces and graph-based methods.
All experiments were evaluated on the IL-PCR test set using the metrics described in Section~3.
Table~\ref{tab:baselines} summarizes the key results across all blocks.

\subsection{Block 1: Raw Text Retrieval}

Three baselines use unprocessed document text directly:
\begin{itemize}
    \item \textbf{BM25} ($k_1{=}1.6$, $b{=}0.75$): Standard probabilistic retrieval.
    \item \textbf{TF-IDF (unigram)}: Cosine similarity on unigram vectors.
    \item \textbf{TF-IDF (uni+bi)}: Cosine similarity on unigram + bigram vectors.
\end{itemize}

TF-IDF consistently outperforms BM25 on this dataset by a wide margin.
The best raw-text variant is \textbf{TF-IDF (uni+bi)}, achieving Micro-F1 of 0.349 at $K{=}10$ and MAP of 0.421.
Adding bigrams provides a modest improvement over unigrams (MAP 0.421 vs.~0.406), suggesting that short legal phrases such as statute references and judicial terminology carry meaningful retrieval signal.
BM25's underperformance is noteworthy: its term frequency saturation parameter $k_1$ is designed for documents of moderate length, but Indian court judgments are extremely long (often exceeding 10,000 tokens), causing the saturation curve to behave suboptimally.

\subsection{Block 2: Standard Preprocessing}

Documents are preprocessed with lowercasing, non-alphabetic removal, lemmatization, and stopword removal.
We evaluate BM25 with three parameter settings, TF-IDF (uni+bi), and LSA with 100 and 200 latent dimensions.

Key findings:
\begin{itemize}
    \item Preprocessing provides minimal benefit over raw text for TF-IDF, confirming that legal vocabulary is sufficiently distinctive.
    \item \textbf{LSA-200} (MRR: 0.693, MAP: 0.405) slightly outperforms LSA-100, confirming that more latent dimensions better capture the semantic variety in legal text, but both fall short of direct TF-IDF.
    \item BM25 with reduced length normalization ($b{=}0.5$) degrades performance significantly.
\end{itemize}

\subsection{Block 3: Concept Extraction}

Documents are reduced to extracted legal concepts using rule-based extraction: citation patterns (e.g., ``2004 SCC 12''), statute references (e.g., ``Section 302''), act names (e.g., ``Indian Penal Code''), and POS-tagged nouns/verbs extracted via spaCy.
TF-IDF on concept-only text achieves MAP of 0.407 and Micro-F1@10 of 0.335, competitive with full-text TF-IDF despite discarding the vast majority of the document.
This is a significant finding: it suggests that legal named entities---statutes, citations, and act references---carry substantial retrieval signal and that much of the surrounding natural language text is redundant for retrieval purposes.
This observation motivates our proposed approach of filtering documents to retain only the most informative segments via rhetorical role prediction.

\subsection{Block 4: BM25 + Graph Ensemble}

A hybrid ensemble combines BM25 scores with entity-graph scores computed via TF-IDF cosine similarity over regex-extracted legal named entities.
Three weighting configurations were tested (BM25:Entity = 80:20, 60:40, 40:60).
All variants underperform pure TF-IDF, with the best (80:20) achieving Micro-F1@10 of only 0.203.
The sparse regex-based entity extraction introduces more noise than signal.

\subsection{Block 5: Graph Diffusion Re-Ranking}

Initial TF-IDF scores are propagated through a candidate-candidate similarity graph built from legal entity overlap, using the diffusion formula:
\begin{equation}
\mathbf{S}_{\text{final}} = (1-\alpha)\,\mathbf{S}_{\text{seed}} + \alpha\,\mathbf{T}\,\mathbf{S}_{\text{seed}}
\end{equation}
where $\mathbf{T}$ is the row-normalized transition matrix.
Three values of $\alpha$ (0.15, 0.35, 0.55) were tested.
Increasing $\alpha$ consistently degrades performance, indicating that the citation graph amplifies errors rather than propagating relevance.

\begin{table*}[t]
\centering
\small
\begin{tabular}{@{}llcccccc@{}}
\toprule
\textbf{Block} & \textbf{Variant} & \textbf{F1@1} & \textbf{F1@5} & \textbf{F1@10} & \textbf{F1@20} & \textbf{MAP@20} & \textbf{MRR@20} \\
\midrule
\multirow{3}{*}{Raw Text}
  & BM25 Standard     & 0.075 & 0.153 & 0.146 & 0.123 & 0.166 & 0.414 \\
  & TF-IDF Unigram    & 0.135 & 0.319 & 0.336 & 0.277 & 0.406 & 0.682 \\
  & TF-IDF Uni+Bi     & 0.134 & \textbf{0.330} & \textbf{0.349} & \textbf{0.285} & \textbf{0.421} & \textbf{0.690} \\
\midrule
\multirow{4}{*}{Std. Prep.}
  & BM25 Base          & 0.078 & 0.144 & 0.136 & 0.117 & 0.163 & 0.415 \\
  & TF-IDF Uni+Bi      & 0.129 & 0.315 & 0.341 & 0.281 & 0.408 & 0.667 \\
  & LSA-100            & 0.125 & 0.298 & 0.313 & 0.267 & 0.374 & 0.656 \\
  & LSA-200            & \textbf{0.139} & 0.317 & 0.330 & 0.275 & 0.405 & 0.693 \\
\midrule
Concept Extr. & TF-IDF & 0.137 & 0.307 & 0.335 & 0.273 & 0.407 & 0.678 \\
\midrule
\multirow{3}{*}{Graph Ens.}
  & BM80:Ent20       & 0.104 & 0.202 & 0.203 & 0.169 & 0.236 & 0.532 \\
  & BM60:Ent40       & 0.095 & 0.197 & 0.197 & 0.170 & 0.217 & 0.497 \\
  & BM40:Ent60       & 0.086 & 0.172 & 0.179 & 0.153 & 0.188 & 0.453 \\
\midrule
\multirow{3}{*}{Graph Diff.}
  & $\alpha{=}0.15$  & 0.121 & 0.281 & 0.296 & 0.248 & 0.346 & 0.624 \\
  & $\alpha{=}0.35$  & 0.117 & 0.273 & 0.284 & 0.247 & 0.336 & 0.607 \\
  & $\alpha{=}0.55$  & 0.111 & 0.256 & 0.274 & 0.237 & 0.311 & 0.583 \\
\bottomrule
\end{tabular}
\caption{Baseline experiment results on the IL-PCR test set. Micro-F1 is reported at cutoffs $K \in \{1, 5, 10, 20\}$. MAP and MRR are computed at $K{=}20$. Best values in each column are \textbf{bolded}.}
\label{tab:baselines}
\end{table*}

%----------------------------------------------------------------------
\section{Hybrid BM25 + Semantic Retrieval}
%----------------------------------------------------------------------

Beyond classical baselines, we evaluated hybrid pipelines that combine BM25 candidate generation with transformer-based semantic re-ranking.
The motivation is that BM25 provides efficient lexical recall while dense encoders capture semantic similarity that lexical methods miss.
All hybrid experiments were conducted on the IL-PCR training split to preserve the test set for final evaluation.
Scores are fused using a linear combination: $s_{\text{final}} = \alpha \cdot s_{\text{BM25}} + (1-\alpha) \cdot s_{\text{semantic}}$ with $\alpha{=}0.3$.
Table~\ref{tab:hybrid} summarizes the results.

\subsection{BM25 Baseline}
BM25 with unigram indexing achieves a best F1 of 11.90\% at $K{=}5$ on the training split.

\subsection{BM25 + InLegalBERT}
We encode documents using \texttt{law-ai/InLegalBERT} with top-512 truncation and fuse scores with BM25 ($\alpha{=}0.3$ BM25, $0.7$ BERT).
This yields a best F1 of only 4.98\% at $K{=}20$---\emph{far worse than BM25 alone}.
Top-512 truncation discards critical content because Indian legal judgments bury relevant material deep in the document; the first 512 tokens are mostly preamble.

\subsection{BM25 + SBERT (Full Document)}
Using \texttt{all-MiniLM-L6-v2} with full-document encoding (internally truncated to 512 tokens) and the same fusion weighting achieves 10.33\% F1 at $K{=}7$---slightly worse than BM25 alone, confirming that naive truncation is ineffective for long legal documents.

\subsection{BM25 + SBERT (Chunked Encoding)}
We encode documents using 256-token chunks with 128-token stride, applying a sliding window across the full document text.
Each chunk is independently encoded by \texttt{all-MiniLM-L6-v2}, and the resulting chunk embeddings are aggregated via max-pooling to form a single document-level representation.
The 128-token overlap ensures that no information is lost at chunk boundaries, and max-pooling selects the most discriminative features across all chunks.
This achieves a \textbf{best F1 of 19.63\%} at $K{=}6$, a \textbf{+7.73 percentage point improvement} over BM25.
This result confirms that chunking is essential for applying semantic models to long legal documents, and that the information lost by truncation-based approaches is substantial and recoverable.

Table~\ref{tab:hybrid_example} shows an example top-10 ranking produced by the chunked hybrid method for a single query.
The two relevant candidates appear at ranks 6 and 8, illustrating that the model assigns non-trivial scores to relevant cases but they are interspersed with irrelevant results.

\begin{table}[htbp]
\centering
\footnotesize
\setlength{\tabcolsep}{3pt}
\begin{tabular}{@{}lccrr@{}}
\toprule
\textbf{Method} & \textbf{Enc.} & \textbf{F1} & \textbf{$K$} & \textbf{$\Delta$} \\
\midrule
BM25              & ---       & 11.90 & 5  & --- \\
+ InLegalBERT     & Top-512   & 4.98  & 20 & $-$6.92 \\
+ SBERT (full)    & 512 tr.   & 10.33 & 7  & $-$1.57 \\
+ SBERT (chunk)   & 256/128   & \textbf{19.63} & 6  & \textbf{+7.73} \\
\bottomrule
\end{tabular}
\caption{Hybrid results on IL-PCR train split (F1 in \%). $\Delta$: pp improvement over BM25. Fusion: $\alpha{=}0.3$ (BM25), $0.7$ (semantic).}
\label{tab:hybrid}
\end{table}

\begin{table}[htbp]
\centering
\footnotesize
\setlength{\tabcolsep}{3pt}
\begin{tabular}{@{}clcc@{}}
\toprule
\textbf{Rank} & \textbf{Retrieved Case} & \textbf{Score} & \textbf{Rel.} \\
\midrule
1  & 0000000599 & 1.0000 & (self) \\
2  & 0001207163 & 0.6317 & \\
3  & 0001296654 & 0.6223 & \\
4  & 0000175790 & 0.5628 & \\
5  & 0000382199 & 0.5424 & \\
6  & 0001742283 & 0.5352 & \checkmark \\
7  & 0001925650 & 0.5346 & \\
8  & 0001292810 & 0.5148 & \checkmark \\
9  & 0001812278 & 0.5128 & \\
10 & 0001831036 & 0.5084 & \\
\bottomrule
\end{tabular}
\caption{Example top-10 ranking for query \texttt{0000000599} using BM25~+~SBERT (256-token chunked). Relevant candidates (\checkmark) appear at ranks 6 and 8.}
\label{tab:hybrid_example}
\end{table}

%----------------------------------------------------------------------
\section{Rhetorical Role Prediction}
%----------------------------------------------------------------------

A key component of our proposed pipeline is rhetorical role segmentation, which decomposes legal judgments into semantically coherent units.
We have deployed a pretrained multi-task learning (MTL) classifier for this purpose.

\subsection{Task Formulation and Metrics}

RR prediction is formulated as a \textbf{sentence-level multi-class classification} task: given a legal document pre-segmented into sentences, the model must assign exactly one rhetorical role label to each sentence.
The IL-TUR benchmark defines 13 labels: \emph{Fact}, \emph{Issue}, \emph{Argument~(Petitioner)}, \emph{Argument~(Respondent)}, \emph{Statute}, \emph{Dissent}, \emph{Precedent Relied Upon}, \emph{Precedent Not Relied Upon}, \emph{Precedent Overruled}, \emph{Ruling By Lower Court}, \emph{Ratio Of The Decision}, \emph{Ruling By Present Court}, and \emph{None}.
The official evaluation metric is \textbf{macro-F1}, which weights all 13 labels equally.

The training dataset \cite{malik2022rr} comprises 21,184 sentences from Indian court judgments on banking and competition law, annotated by six legal experts from a reputed law school.
Each sentence includes primary, secondary, and tertiary annotations per expert, with a final consolidated label.
Because legal documents exhibit strong label inertia---adjacent sentences typically share the same role---sequence-labeling architectures (BiLSTM-CRF) substantially outperform independent classifiers.

Table~\ref{tab:rr_baselines} compares the IL-TUR baseline models.
The MTL-BiLSTM-CRF model achieves the best macro-F1 of 0.70/0.69, confirming that the auxiliary shift-prediction task helps the main classifier detect label transitions.

\begin{table}[htbp]
\centering
\small
\begin{tabular}{@{}lccc@{}}
\toprule
\textbf{Model} & \textbf{P} & \textbf{R} & \textbf{F1} \\
\midrule
BiLSTM-CRF (sent2vec)        & 0.61 & 0.59 & 0.65 \\
BiLSTM-CRF (BERT embs)       & 0.63 & 0.63 & 0.63 \\
LSP-BiLSTM-CRF (BERT-SC)     & 0.68 & 0.65 & 0.67 \\
MTL-BiLSTM-CRF (BERT-SC)     & \textbf{0.69} & \textbf{0.70} & \textbf{0.70} \\
\bottomrule
\end{tabular}
\caption{RR prediction baselines from \citet{malik2022rr} (macro-averaged). The MTL model uses shift prediction as an auxiliary task.}
\label{tab:rr_baselines}
\end{table}

\subsection{Model Architecture}

The model combines BERT embeddings with stacked BiLSTM layers and a CRF decoder:

\begin{enumerate}
    \item \textbf{Sentence Embedding}: Each sentence is encoded using \texttt{bert-base-uncased}; the \texttt{[CLS]} embedding (768-dim) serves as the sentence representation.
    \item \textbf{Binary/Shift Module}: A BiLSTM operates on context-enriched embeddings formed by concatenating the previous, current, and next sentence embeddings (2304-dim), learning coarse structural shift signals.
    \item \textbf{Main Classifier}: A second BiLSTM takes the 768-dim BERT embeddings concatenated with hidden states from the binary module (1536-dim total), enabling cross-module fusion.
    \item \textbf{CRF Decoding}: Conditional Random Field decoders model label transitions across sentences, producing globally coherent label sequences via Viterbi decoding.
\end{enumerate}

\subsection{Predicted Labels}

Although the IL-TUR benchmark defines 13 labels, the pretrained model we deploy uses a reduced label set of six roles: \textbf{Fact}, \textbf{Precedent}, \textbf{RulingByLowerCourt}, \textbf{Argument}, \textbf{Ratio Decidendi}, and \textbf{Statute}.
The reduced set collapses related categories (e.g., the three precedent sub-types into a single \emph{Precedent} label) and omits infrequent labels such as \emph{Dissent} and \emph{Issue}.
In a test on a 1963 Indian Supreme Court judgment (185 sentences), the model produces coherent predictions: \emph{Fact} labels dominate the early sections, \emph{Precedent} labels appear in the analytical middle, and \emph{Ruling} labels cluster in the preamble and conclusion.
Table~\ref{tab:rr_example} illustrates representative predictions from this judgment.

\begin{table}[htbp]
\centering
\footnotesize
\setlength{\tabcolsep}{4pt}
\begin{tabular}{@{}p{0.62\columnwidth}l@{}}
\toprule
\textbf{Sentence (excerpt)} & \textbf{Predicted Role} \\
\midrule
``The validity of the election of the appellant to [ORG] at the third general elections held in the month of February, 1962, was challenged by two of the electors of the constituency\ldots'' & Fact \\
\midrule
``The foundation of the argument is that there has been a non-compliance with the provisions of s.~82.'' & Precedent \\
\midrule
``Appeals by special leave from the judgment and order dated August 31, 1962, of [ORG] in D. [NAME] Civil Writ Petitions Nos.~376 and 377 of 1962.'' & RulingByLowerCourt \\
\bottomrule
\end{tabular}
\caption{Example rhetorical role predictions on a 1963 Indian Supreme Court judgment. Named entities are anonymised as \texttt{[ORG]} and \texttt{[NAME]}.}
\label{tab:rr_example}
\end{table}

To assess inference stability, we ran the classifier three times on a 228-sentence judgment.
Table~\ref{tab:rr_ndet} summarises the non-determinism analysis.
All disagreements are confined to a contiguous block of sentences (4--118) in Run~1, which diverges into \texttt{Precedent} while Runs~2 and~3 consistently predict \texttt{RulingByLowerCourt}.
This is caused by floating-point variance on CPU cascading through the CRF decoder.
For production use, majority voting or fixed random seeds can mitigate this.

\begin{table}[htbp]
\centering
\small
\begin{tabular}{@{}lr@{}}
\toprule
\textbf{Metric} & \textbf{Value} \\
\midrule
Total sentences          & 228 \\
Full agreement (3/3 runs) & 113 (49.6\%) \\
Disagreement (Run~1 $\neq$ Runs~2,3) & 115 (50.4\%) \\
Runs~2 vs.~3 agreement   & 228/228 (100\%) \\
\bottomrule
\end{tabular}
\caption{Non-determinism analysis of the rhetorical role classifier across three inference runs on a 228-sentence Indian Supreme Court judgment. Disagreements are confined to sentences 4--118.}
\label{tab:rr_ndet}
\end{table}

\subsection{Integration with Retrieval}

For the proposed retrieval pipeline, we extract only the \textbf{Facts} and \textbf{Ratio Decidendi} segments from each document.
These filtered texts are then indexed and used for retrieval, reducing noise from procedural details, arguments, and other sections that do not directly contribute to precedent similarity.

%----------------------------------------------------------------------
\section{Legal Statute Identification}
%----------------------------------------------------------------------

As a complementary component, we evaluated a pretrained Legal Statute Identification (LSI) model that predicts applicable Indian Penal Code (IPC) sections from case text.
The IPC is the principal criminal code of India, covering offences ranging from theft and assault to murder and treason.
Identifying which IPC sections apply to a given case is a critical step in the Indian judicial process, traditionally performed manually by police personnel and legal professionals.
Automating this step not only reduces workload but also provides structured statute-level features that can be exploited for graph-based retrieval.

\subsection{Model Architecture}

The model uses a \textbf{Hierarchical BERT (HierBERT)} architecture, which is specifically designed to handle documents that exceed the 512-token limit of standard BERT encoders.
Since fact descriptions in the ILSI dataset are pre-segmented into sentences, the hierarchical approach processes one sentence at a time:
\begin{enumerate}
    \item \textbf{Sentence Encoding}: Each sentence is individually passed through \texttt{InLegalBERT} \cite{paul2022ilsi}, a BERT model pretrained on Indian legal corpora; the \texttt{[CLS]} embedding is extracted as the sentence representation.
    \item \textbf{Document Encoding}: The sequence of \texttt{[CLS]} embeddings is fed into a Bi-LSTM with attention, which aggregates sentence-level representations into a single document-level vector. The attention mechanism allows the model to focus on the most informative sentences for statute prediction.
    \item \textbf{Classification}: A fully connected layer projects the document vector to a 100-dimensional output space (corresponding to the 100 most frequent IPC sections) with sigmoid activation for multi-label prediction. Sections with predicted probability $> 0.5$ are output as relevant.
\end{enumerate}

\subsection{Task Formulation and Metrics}

LSI is formulated as a \textbf{multi-label text classification} task: given the fact portion of a case document, the model must predict one or more applicable statutes from the 100 most frequent IPC sections.
The ILSI dataset \cite{paul2022ilsi} consists of $\sim$65k criminal court documents from the Supreme Court of India and six High Courts, with anonymised named entities.
Labels with predicted probability $> 0.5$ are considered relevant.

The official IL-TUR evaluation metric is \textbf{macro-averaged F1}, which weights all 100 statute classes equally regardless of frequency.
This is complemented by macro-averaged precision and recall.
In the IL-TUR leaderboard (Table~\ref{tab:iltur}), InLegalBERT achieves the highest LSI score (26.23 macro-F1), followed by CaseLawBERT and LegalBERT, indicating that domain-specific pretraining is beneficial for statute identification---unlike the PCR task where generic BERT performs better.

\subsection{Results}

We evaluated the pretrained model on four synthetic case summaries spanning distinct categories of IPC offences.
Table~\ref{tab:lsi} summarises the predictions.
The model produces \textbf{no false positives} across all four samples and correctly identifies the core applicable statutes with high confidence.
Missed sections (e.g., Section~34 at 0.347, Section~376(2) at 0.490) consistently appear just below the 0.5 threshold.
The cheating/forgery case achieves a perfect 5/5 prediction.
This model can be used to extract statute-level features for graph-based retrieval components.

\begin{table}[htbp]
\centering
\footnotesize
\setlength{\tabcolsep}{3pt}
\begin{tabular}{@{}llcl@{}}
\toprule
\textbf{Case Type} & \textbf{Expected} & \textbf{Correct} & \textbf{Missed} \\
\midrule
Murder / Dowry   & 302, 304B, 498A, 34    & 3/4 & 34 (0.347) \\
Robbery + Hurt   & 392, 394, 397, 120B    & 3/4 & 120B \\
Cheating / Forgery & 420, 467, 468, 471, 415 & 5/5 & --- \\
Kidnapping / Assault & 363, 366, 366A, 376, 376(2) & 4/5 & 376(2) (0.490) \\
\bottomrule
\end{tabular}
\caption{LSI model predictions on four synthetic case summaries (threshold~=~0.5). \emph{Correct} shows the number of expected IPC sections predicted above threshold. Parenthetical values are near-miss confidences.}
\label{tab:lsi}
\end{table}

\begin{table*}[htbp]
\centering
\small
\begin{tabular}{@{}llccc@{}}
\toprule
\textbf{Method} & \textbf{Submitted By} & \textbf{RR} & \textbf{LSI} & \textbf{PCR} \\
\midrule
SOTA             & multiple & 69.01 & 28.08 & 39.15 \\
BERT             & multiple & 58    & 18.44 & 9.24  \\
LegalBERT        & multiple & 54    & 21.74 & 8.67  \\
InLegalBERT      & multiple & 58    & 26.23 & 7.57  \\
\midrule
GPT-3.5 (0-shot) & IL-TUR   & 30.95 & 21.55 & ---   \\
GPT-3.5 (1-shot) & IL-TUR   & 30.05 & 22.61 & ---   \\
GPT-3.5 (2-shot) & IL-TUR   & 30.31 & 21.40 & ---   \\
GPT-4 (0-shot)   & IL-TUR   & 37.37 & 23.99 & ---   \\
GPT-4 (1-shot)   & IL-TUR   & 37.43 & 22.26 & ---   \\
GPT-4 (2-shot)   & IL-TUR   & 38.18 & 20.53 & ---   \\
\midrule
Ideal Baseline   & dummy    & 100.00 & 100.00 & 64.87 \\
Random Baseline  & dummy    & 10.70  & 4.78   & 0.49  \\
\bottomrule
\end{tabular}
\caption{IL-TUR leaderboard scores for Rhetorical Role (RR), Legal Statute Identification (LSI), and Prior Case Retrieval (PCR) tasks. BERT outperforms domain-adapted models (LegalBERT, InLegalBERT) on PCR despite lacking legal pretraining.}
\label{tab:iltur}
\end{table*}

%----------------------------------------------------------------------
\section{Key Findings So Far}
%----------------------------------------------------------------------

\begin{enumerate}
    \item \textbf{TF-IDF outperforms BM25} on long legal documents by a substantial margin (MAP 0.421 vs.~0.166). BM25's term saturation assumptions, which cap the contribution of repeated terms via the $k_1$ parameter, appear poorly suited to the very long documents in legal corpora where important terms (e.g., specific statute references) may appear dozens of times. TF-IDF's unbounded term weighting better captures these high-frequency discriminative terms.

    \item \textbf{Raw text is competitive with preprocessed text} for TF-IDF retrieval. Standard NLP preprocessing (lowercasing, lemmatization, stopword removal) provides negligible improvement, suggesting that legal vocabulary is sufficiently distinctive and domain-specific that these generic techniques add minimal value. This is consistent with findings in other specialized domains such as biomedical text retrieval.

    \item \textbf{Domain-specific pretraining hurts PCR.} As per the IL-TUR leaderboard (Table~\ref{tab:iltur}), generic BERT (micro-F1 9.24) outperforms LegalBERT (8.67) and InLegalBERT (7.57) on the PCR task. This counter-intuitive result suggests that domain-adapted models may overfit to the distribution of legal language at the expense of the more abstract reasoning required for precedent retrieval. However, the trend reverses for LSI, where InLegalBERT achieves the best score (26.23), indicating that domain pretraining helps for classification tasks but not necessarily for retrieval.

    \item \textbf{Chunking is essential for semantic models.} SBERT with 256-token chunks and 128-token stride achieves a +7.73pp gain over BM25, while full-document encoding (truncated to 512 tokens) and top-512 approaches actually \emph{hurt} performance. This demonstrates that the information content of legal documents is distributed throughout the text, and discarding everything beyond the first 512 tokens loses critical factual and reasoning content.

    \item \textbf{Concept extraction is surprisingly effective.} TF-IDF over regex-extracted legal concepts (citations, statutes, act names) achieves MAP 0.407---only 3\% below full-text TF-IDF despite discarding the majority of the document text. This validates our hypothesis that filtering documents to their most informative content can maintain or improve retrieval quality.
\end{enumerate}

%----------------------------------------------------------------------
\section{Remaining Work}
%----------------------------------------------------------------------

The following tasks remain for the second half of the project:

\begin{enumerate}
    \item \textbf{Rhetorical role--filtered retrieval}: Apply the RR segmentation model to the full IL-PCR corpus, extract Facts and Ratio segments, and perform BM25 and TF-IDF retrieval on the filtered text. We hypothesize this will reduce noise and improve precision.

    \item \textbf{Segment-aware semantic retrieval}: Apply chunked SBERT encoding specifically to the extracted Facts and Ratio segments, combining rhetorical filtering with the best-performing semantic method.

    \item \textbf{Statute-based graph features}: Use the LSI model to extract statute predictions for each document and construct a case--statute bipartite graph. Cases sharing the same statutes will receive a graph-based similarity boost.

    \item \textbf{Full test-set evaluation}: Run all methods on the IL-PCR test split and report final comparative results.

    \item \textbf{Ablation studies}: Evaluate the contribution of each pipeline component (RR filtering, chunked SBERT, statute graph) in isolation and in combination.

    \item \textbf{Optional SBERT re-ranking}: Implement a two-stage pipeline where BM25 generates top-100 candidates and chunked SBERT re-ranks them, potentially combined with RR filtering.
    
    \item \textbf{Expanding RR labels}: We will evaluate the performance upon extracting more data from the RR predictions, not just facts and ratio decidendi.
\end{enumerate}

\section{Timeline}
\begin{itemize}
    \item \textbf{Week 1}: Apply rhetorical role segmentation to the full IL-PCR corpus and extract \textbf{Facts} and \textbf{Ratio Decidendi} sections from each document. Implement \textbf{BM25 and TF-IDF retrieval} using the RR-filtered text and compare performance with raw-text baselines.
    \item \textbf{Week 2}: Implement \textbf{segment-aware semantic retrieval} using \textbf{chunked SBERT embeddings} on extracted sections.
    \item \textbf{Week 3}: Build \textbf{statute-based graph features} using predictions from the Legal Statute Identification model and integrate them into the retrieval pipeline
    \item \textbf{Week 4}: Conduction ablation studies and a full evaluation on the IL-PCR test set, comparing all methods and reporting final results.
\end{itemize}

%----------------------------------------------------------------------
\section{Conclusion}
%----------------------------------------------------------------------

We have completed extensive baseline experiments for the prior case retrieval task on the IL-PCR benchmark, evaluating classical lexical methods, latent semantic approaches, concept extraction, graph-based ensembles, graph diffusion, and hybrid semantic pipelines.
Our results establish that TF-IDF with bigrams is the strongest classical baseline (MAP 0.421), and that hybrid BM25+SBERT with chunked encoding substantially outperforms all classical methods (F1 19.63\%, a +7.73pp gain).
We have also demonstrated working rhetorical role prediction and legal statute identification components: the RR classifier produces coherent sentence-level labels across long judgments, and the LSI model achieves zero false positives on synthetic test cases.

Two cross-cutting themes emerge from our experiments.
First, long document handling is the central challenge in legal retrieval---methods that truncate or ignore document structure consistently underperform those that process the full text.
Second, content filtering (whether via concept extraction or rhetorical role segmentation) is a promising strategy: even aggressive filtering to legal entities preserves most retrieval signal.

The remaining work focuses on integrating these components into a unified retrieval pipeline that filters documents by rhetorical role before applying semantic retrieval.
We expect that restricting the index to \emph{Facts} and \emph{Ratio Decidendi} segments will reduce noise from procedural details and arguments, improving both precision and recall.
Additionally, statute-based graph features from the LSI model will provide a complementary retrieval signal based on shared legal provisions between cases.

\section{Links}
\begin{itemize}
    \item \textbf{Code repository}: \url{https://github.com/gautambhetanabhotla/inlp-project}
    \item \textbf{Dataset}: \url{https://drive.google.com/drive/folders/1UoATAoSnTboJMztTN-EPHKorkYHpOV9x?usp=sharing}
\end{itemize}

\bibliography{custom}

\end{document}